%!TeX root=../emmatop.tex
\chapter[Chapter \thechapter]{}
\lettrine[lraise=0.3]{M}{r} Weston was a native of Highbury, and born of a respectable family, which for the last two or three generations had been rising into gentility and property. He had received a good education, but, on succeeding early in life to a small independence, had become indisposed for any of the more homely pursuits in which his brothers were engaged, and had satisfied an active, cheerful mind and social temper by entering into the militia of his county, then embodied.

Captain Weston was a general favourite; and when the chances of his military life had introduced him to Miss Churchill, of a great Yorkshire family, and Miss Churchill fell in love with him, nobody was surprized, except her brother and his wife, who had never seen him, and who were full of pride and importance, which the connexion would offend.

Miss Churchill, however, being of age, and with the full command of her fortune—though her fortune bore no proportion to the family-estate—was not to be dissuaded from the marriage, and it took place, to the infinite mortification of Mr and Mrs Churchill, who threw her off with due decorum. It was an unsuitable connexion, and did not produce much happiness. Mrs Weston ought to have found more in it, for she had a husband whose warm heart and sweet temper made him think every thing due to her in return for the great goodness of being in love with him; but though she had one sort of spirit, she had not the best. She had resolution enough to pursue her own will in spite of her brother, but not enough to refrain from unreasonable regrets at that brother's unreasonable anger, nor from missing the luxuries of her former home. They lived beyond their income, but still it was nothing in comparison of Enscombe: she did not cease to love her husband, but she wanted at once to be the wife of Captain Weston, and Miss Churchill of Enscombe.

Captain Weston, who had been considered, especially by the Churchills, as making such an amazing match, was proved to have much the worst of the bargain; for when his wife died, after a three years' marriage, he was rather a poorer man than at first, and with a child to maintain. From the expense of the child, however, he was soon relieved. The boy had, with the additional softening claim of a lingering illness of his mother's, been the means of a sort of reconciliation; and Mr and Mrs Churchill, having no children of their own, nor any other young creature of equal kindred to care for, offered to take the whole charge of the little Frank soon after her decease. Some scruples and some reluctance the widower-father may be supposed to have felt; but as they were overcome by other considerations, the child was given up to the care and the wealth of the Churchills, and he had only his own comfort to seek, and his own situation to improve as he could.

A complete change of life became desirable. He quitted the militia and engaged in trade, having brothers already established in a good way in London, which afforded him a favourable opening. It was a concern which brought just employment enough. He had still a small house in Highbury, where most of his leisure days were spent; and between useful occupation and the pleasures of society, the next eighteen or twenty years of his life passed cheerfully away. He had, by that time, realised an easy competence—enough to secure the purchase of a little estate adjoining Highbury, which he had always longed for—enough to marry a woman as portionless even as Miss Taylor, and to live according to the wishes of his own friendly and social disposition.

It was now some time since Miss Taylor had begun to influence his schemes; but as it was not the tyrannic influence of youth on youth, it had not shaken his determination of never settling till he could purchase Randalls, and the sale of Randalls was long looked forward to; but he had gone steadily on, with these objects in view, till they were accomplished. He had made his fortune, bought his house, and obtained his wife; and was beginning a new period of existence, with every probability of greater happiness than in any yet passed through. He had never been an unhappy man; his own temper had secured him from that, even in his first marriage; but his second must shew him how delightful a well-judging and truly amiable woman could be, and must give him the pleasantest proof of its being a great deal better to choose than to be chosen, to excite gratitude than to feel it.

He had only himself to please in his choice: his fortune was his own; for as to Frank, it was more than being tacitly brought up as his uncle's heir, it had become so avowed an adoption as to have him assume the name of Churchill on coming of age. It was most unlikely, therefore, that he should ever want his father's assistance. His father had no apprehension of it. The aunt was a capricious woman, and governed her husband entirely; but it was not in Mr Weston's nature to imagine that any caprice could be strong enough to affect one so dear, and, as he believed, so deservedly dear. He saw his son every year in London, and was proud of him; and his fond report of him as a very fine young man had made Highbury feel a sort of pride in him too. He was looked on as sufficiently belonging to the place to make his merits and prospects a kind of common concern.

Mr Frank Churchill was one of the boasts of Highbury, and a lively curiosity to see him prevailed, though the compliment was so little returned that he had never been there in his life. His coming to visit his father had been often talked of but never achieved.

Now, upon his father's marriage, it was very generally proposed, as a most proper attention, that the visit should take place. There was not a dissentient voice on the subject, either when Mrs Perry drank tea with Mrs and Miss Bates, or when Mrs and Miss Bates returned the visit. Now was the time for Mr Frank Churchill to come among them; and the hope strengthened when it was understood that he had written to his new mother on the occasion. For a few days, every morning visit in Highbury included some mention of the handsome letter Mrs Weston had received. <I suppose you have heard of the handsome letter Mr Frank Churchill has written to Mrs Weston? I understand it was a very handsome letter, indeed. Mr Woodhouse told me of it. Mr Woodhouse saw the letter, and he says he never saw such a handsome letter in his life.>

It was, indeed, a highly prized letter. Mrs Weston had, of course, formed a very favourable idea of the young man; and such a pleasing attention was an irresistible proof of his great good sense, and a most welcome addition to every source and every expression of congratulation which her marriage had already secured. She felt herself a most fortunate woman; and she had lived long enough to know how fortunate she might well be thought, where the only regret was for a partial separation from friends whose friendship for her had never cooled, and who could ill bear to part with her.

She knew that at times she must be missed; and could not think, without pain, of Emma's losing a single pleasure, or suffering an hour's ennui, from the want of her companionableness: but dear Emma was of no feeble character; she was more equal to her situation than most girls would have been, and had sense, and energy, and spirits that might be hoped would bear her well and happily through its little difficulties and privations. And then there was such comfort in the very easy distance of Randalls from Hartfield, so convenient for even solitary female walking, and in Mr Weston's disposition and circumstances, which would make the approaching season no hindrance to their spending half the evenings in the week together.

\begin{figure}[tbph]
\centering
\includegraphics[width=.8\linewidth]{2cake}
\caption{With a slice of Mrs Weston's wedding-cake}
\end{figure}

Her situation was altogether the subject of hours of gratitude to Mrs Weston, and of moments only of regret; and her satisfaction—her more than satisfaction—her cheerful enjoyment, was so just and so apparent, that Emma, well as she knew her father, was sometimes taken by surprize at his being still able to pity <poor Miss Taylor,> when they left her at Randalls in the centre of every domestic comfort, or saw her go away in the evening attended by her pleasant husband to a carriage of her own. But never did she go without Mr Woodhouse's giving a gentle sigh, and saying, <Ah, poor Miss Taylor! She would be very glad to stay.>

There was no recovering Miss Taylor—nor much likelihood of ceasing to pity her; but a few weeks brought some alleviation to Mr Woodhouse. The compliments of his neighbours were over; he was no longer teased by being wished joy of so sorrowful an event; and the wedding-cake, which had been a great distress to him, was all eat up. His own stomach could bear nothing rich, and he could never believe other people to be different from himself. What was unwholesome to him he regarded as unfit for any body; and he had, therefore, earnestly tried to dissuade them from having any wedding-cake at all, and when that proved vain, as earnestly tried to prevent any body's eating it. He had been at the pains of consulting Mr Perry, the apothecary, on the subject. Mr Perry was an intelligent, gentlemanlike man, whose frequent visits were one of the comforts of Mr Woodhouse's life; and upon being applied to, he could not but acknowledge (though it seemed rather against the bias of inclination) that wedding-cake might certainly disagree with many—perhaps with most people, unless taken moderately. With such an opinion, in confirmation of his own, Mr Woodhouse hoped to influence every visitor of the newly married pair; but still the cake was eaten; and there was no rest for his benevolent nerves till it was all gone.




There was a strange rumour in Highbury of all the little Perrys being seen with a slice of Mrs Weston's wedding-cake in their hands: but Mr Woodhouse would never believe it.